% sections/1_ref_info.tex
\chapter{Справочные сведения о походе}
\label{ch:refinfo}

% ---------------- 1.1 Общие сведения ----------------
\section{Общие сведения о походе}
\begin{tourinfotable}
\tourinfo{Район похода}{Центральный Тянь-Шань, Западная часть хребта; Терскей-Алатоо}
\tourinfo{Вид туризма}{Горный}
\tourinfo{Категория сложности похода}{Третья}
\tourinfo{Проводящая организация}{Горная секция МФТИ, г. Долгопрудный}
\tourinfo{Руководитель}{Вишнякова Анастасия Васильевна}
\tourinfo{Контакты руководителя}{e-mail: \texttt{luna-mars@bk.ru}}
\tourinfo{Сроки (активная часть и полные)}{с 16 июля по 1 августа 2022 года (активная часть)}
\tourinfo{Протяженность по карте (с коэффициентом 1{,}2)}{119 км (142{,}8 км)}
\tourinfo{Протяженность по GPS-треку (из них в зачёт)}{133 км (из них в зачёт — 126{,}5 км)}
\tourinfo{Продолжительность активной части}{15 дней (15 ходовых)}
\tourinfo{Количество днёвок}{0}
\tourinfo{Количество полуднёвок}{0}
\tourinfo{Суммарный набор / сброс высоты}{\,+7377 м / \,-6943 м}
\tourinfo{Максимальная высота (переход)}{4421 м (перелаз между пер. Каракоман Ц. и Джентельмен)}
\tourinfo{Максимальная высота ночёвки}{4249 м}
\end{tourinfotable}

% ---------------- 1.2 Запланированная нитка маршрута ----------------
\section{Запланированная нитка маршрута}
\begin{routeblock}
Пос. Турасу — дол. р. Тон — дол. р. Зындан —
[пер. 267 (1А–1Б, 4000 п/п) — пер. Келъторташ (1Б, 4117)] —
пер. Борду (нк, 3670) —
дол. р. Самынсу —
пер. 255 (1Б, 4200 п/п) —
дол. р. Джеруй —
пер. Джентельмен (2А, 4100)
\end{routeblock}

% ---------------- 1.3 Пройденная нитка маршрута ----------------
\section{Пройденная нитка маршрута}
\begin{routeblock}
Пос. Турасу — дол. р. Тон — дол. р. Зындан —
пер. 267 —
пер. Келъторташ —
...
пер. Джентельмен
\end{routeblock}

% ---------------- 1.4 Определяющие препятствия ----------------
\section{Определяющие препятствия}
\begin{obstacletable}
\obstaclerow{пер. 267}{1Б}{Подход под перевальный «Недооцененный»; взлет 4:44 ЧХВ. (1Б, 4086 п/п)}{Подход по крупной и средней осыпи 15--20\textdegree}
% дополнительные строки по образцу:
% \obstaclerow{пер. 255 Самын-су}{1Б}{Описание времени/перепада}{Характер склонов}
\end{obstacletable}

% ---------------- 1.5 Состав группы ----------------
\section{Состав группы}
\addparticipant{1}{Вишнякова Анастасия Васильевна}{Долгопрудный}{Руководитель}
\addparticipant{2}{Иванов Иван Иванович}{Москва}{Участник}
\addparticipant{3}{Петров Петр Петрович}{Санкт-Петербург}{Участник}

\printparticipants